% mode: ini-etex
% Aimed at pseudo_start, pseudo_input and pseudo_close, and scan_general_text with each of its callers' token
% builders: §1754's cutting into lines, §1755's three ways to print "( ", §1756's trimming, how a pseudo file
% ends (\everyeof, \endinput, errors in it, \end), and what scan_general_text does with a missing brace and
% an \outer macro in its text: the runaway must show its text, which is on the builder.
\input prelude
\scrollmode
\def\space{ }
\newlinechar=`\^^J
\catcode`\^^L=13 \outer\def^^L{\par}
\scantokens{\immediate\write16{one}}
\scantokens{\immediate\write16{two}\immediate\write16{\the\inputlineno}^^J\immediate\write16{\the\inputlineno}}
\immediate\write16{\the\inputlineno}
\scantokens{}
\scantokens{ }
\scantokens{^^J}
\scantokens{^^J^^J}
\scantokens\expandafter{\space}
\def\x{\immediate\write16{from x}}\scantokens\expandafter{\x}
% it reads what show_token_list prints: a control word gets a space, and a # is doubled
\everyeof{\noexpand}\endlinechar=-1
\def\sh#1{\edef\y{\scantokens{#1}}\immediate\write16{[\detokenize\expandafter{\y}]}}
\sh{abc}\sh{\relax x}\sh{\relax\relax}\sh{\.\.}\sh{##}\sh{a b  c}\sh{\ }\sh{{}}\sh{{{a}}}
\catcode`\~=13 \def~{tilde}\sh{~}\sh{\string~}\sh{\noexpand~}
\catcode`\%=12 \sh{a%b}\catcode`\%=14
\sh{a^^Jb}\sh{a^^J^^Jb}\sh{a  ^^J  b}\sh{^^Ja}\sh{a^^J}
\newlinechar=-1 \sh{a^^Jb}\newlinechar=`\^^J
\endlinechar=`\| \sh{a^^Jb}\sh{a}\sh{}\sh{^^J}\endlinechar=-1
\catcode`\ =12 \sh{  a  }\sh{a    }\sh{    }\catcode`\ =10
\everyeof{}\endlinechar=13
% the famous one: the end of the file while a definition is open
\edef\y{\scantokens{a}}\show\y
\def\z{\scantokens{\def\w}}\z{b}
\show\w
\scantokens{\def\w}{c}\show\w
\scantokens{\def\w{d}}\show\w
\scantokens{\expandafter\def\csname v\endcsname{e}}\show\v
\scantokens{\begingroup}\endgroup
\scantokens{\iftrue}\fi
\scantokens{\iffalse}\fi
\scantokens{\iffalse}\fi
\scantokens{\iftrue\else}\fi
\scantokens{\ifcase 1\or}\fi
% an error in it shows its line
\scantokens{a^^Jb\undefined c^^Jd}\par
\tracingscantokens=1
\scantokens{\undefined}
\scantokens{\immediate\write16{in}}
\immediate\write16{}\scantokens{}
\message{some text}\scantokens{}
\message{some text that is getting to be a long line for the log file, nearly fullxxxx}\scantokens{}
\message{some text that is getting to be a long line for the log file, nearly full}\scantokens{}
\immediate\write16{}\scantokens{\scantokens{\scantokens{\undefined}}}
\immediate\write16{}\scantokens{^^J\undefined^^J}
\tracingscantokens=0
% levels: an \input and \endinput inside, nesting, and \inputlineno
\def\in{\immediate\write16{\the\inputlineno}}
\scantokens{\in^^J\scantokens{\in^^J\in}\in^^J\in}
\scantokens{\in\endinput\in^^J\in}\in
\scantokens{\scantokens{\in\endinput}\in}
\everyeof{\immediate\write16{eof}}\scantokens{\in}\scantokens{}\scantokens{\scantokens{}}\scantokens{\endinput}
\everyeof{\endinput}\scantokens{\in}
\everyeof{\scantokens{}}\everyeof{}
\everyeof{\everyeof{}\scantokens{\in}}\scantokens{}
\everyeof{}
% expansion context
\ifnum\scantokens{1}<2 \immediate\write16{less}\fi
\count1=\scantokens{5} \showthe\count1
\expandafter\show\scantokens{\relax}
\immediate\write16{\the\inputlineno}
\show\scantokens
\begingroup\tracingcommands=2 \tracingmacros=2 \def\a{b}\scantokens{\a}\endgroup\par
% scan_left_brace: spaces and \relax are skipped, an implicit brace is taken, others are not
\scantokens \relax {\in}
\let\bg={ \scantokens\bg\in}
\scantokens x\in}
\par
\scantokens\undefined{\in}
\par
\scantokens$\in}
\par
\scantokens}\in}
\par
\immediate\write16{tail}
% the outer macro: a runaway text, with what was scanned so far
\scantokens{abc^^L def}
\par
\scantokens{abc
\outer\def\o{}
\par
\scantokens{abc\o def}
\par
\scantokens\o{abc}
\par
\edef\y{x\scantokens\o{abc}}
\par
\edef\y{x\scantokens{abc\o def}}
\par
\immediate\write16{\detokenize{abc\o def}}
\par
\showtokens{abc\o def}
\par
\def\wrap{\scantokens}\wrap{abc\o def}
\par
\scantokens{a\scantokens{b\o c}d}
\par
\immediate\write16{tail}
% unbalanced at the end of an inner file: the text runs across files
\scantokens{\scantokens{a}\in}
\par
\scantokens{\expandafter\scantokens\expandafter{\the\inputlineno\in}}
\end
